% Generated draft: requires mathematical and source review.
\documentclass[11pt]{article}
\usepackage[margin=25mm]{geometry}
\usepackage{fontspec}
\setmainfont{Latin Modern Roman}
\usepackage{amsmath,amssymb,mathtools,unicode-math}
\setmathfont{Latin Modern Math}
\usepackage{enumitem,needspace,xurl,hyperref}
\setlength{\parindent}{0pt}
\setlength{\parskip}{5pt}
\setlength{\emergencystretch}{3em}
\newcommand{\N}{\mathbb{N}}
\newcommand{\Z}{\mathbb{Z}}
\newcommand{\Q}{\mathbb{Q}}
\newcommand{\R}{\mathbb{R}}
\newcommand{\C}{\mathbb{C}}
\newcommand{\F}{\mathbb{F}}
\newcommand{\A}{\mathbb{A}}
\newcommand{\PP}{\mathbb{P}}
\newcommand{\E}{\mathbb{E}}
\newcommand{\eps}{\varepsilon}
\newcommand{\dd}{\,\mathrm{d}}
\begin{document}
\section*{UnsolvedMath ID: 20003108}
\textbf{Source problem code:} catalogue entry
\subsection{Definitions and mathematical statement}
Fix a finite-type hyperbolic-surface Teichm\"uller space on which the intended Thurston distance is
\[
d_{\mathrm{Th}}(Y,X)=\log\sup_{\lambda\ne0}\frac{\ell_\lambda(X)}{\ell_\lambda(Y)},
\]
with the appropriate measured-lamination convention. Put $d^{\mathrm{rev}}(X,Y)=d_{\mathrm{Th}}(Y,X)$ and fix a basepoint $b$. The reverse horofunction embedding is
\[
Y\longmapsto h_Y^{\mathrm{rev}},\qquad
h_Y^{\mathrm{rev}}(X)=d_{\mathrm{Th}}(Y,X)-d_{\mathrm{Th}}(Y,b).
\]
Its boundary consists of limits in the topology of uniform convergence on compact subsets, excluding functions arising from interior points. Walsh's source question is to identify this reverse horofunction boundary, not the forward boundary. The extracted following paragraph defines reduced Teichm\"uller space for bordered surfaces, but does not clearly state whether it belongs to this problem; bordered/arc-metric conventions must be settled separately.

The aid's partial theorem concerns an unoriented essential nonperipheral simple closed curve $\alpha$, its Dehn twist $T_\alpha$, and any initial point $p$. It asserts, as $|n|\to\infty$,
\[
\exp d_{\mathrm{Th}}(T_\alpha^n p,X)
=|n|A_{p,\alpha}\ell_\alpha(X)+o(|n|),\qquad A_{p,\alpha}>0,
\]
locally uniformly in $X$. In the length-supremum convention its coefficient is the corresponding finite positive supremum of $i(\alpha,\lambda)/\ell_\lambda(p)$; positivity presupposes a lamination intersecting $\alpha$. Normalizing at $b$ cancels the coefficient and yields
\[
h_{T_\alpha^n p}^{\mathrm{rev}}(X)\longrightarrow
h_\alpha(X)=\log\frac{\ell_\alpha(X)}{\ell_\alpha(b)}.
\]
The limit is independent of $p$ and of the sign of the twisting direction. It is a boundary function: pinching $\alpha$ makes it unbounded below, whereas every fixed-point function $h_Y^{\mathrm{rev}}$ is bounded below by $-d_{\mathrm{Th}}(Y,b)$.

The aid further asserts that the assignment $\alpha\mapsto h_\alpha$ is injective on unoriented essential curve isotopy classes. Equivariance uses the basepoint-normalized action
\[
(f\cdot h)(X)=h(f^{-1}X)-h(f^{-1}b),\qquad f\cdot h_\alpha=h_{f\alpha}.
\]
It is not plain precomposition when $f$ moves $b$. This identifies a family of explicit reverse boundary points, not every point of the boundary.

According to the aid the functional form matches Walsh's proposed Busemann formula for a one-component simple stump. It does not show that these points are Busemann points, meaning realizable by the appropriate almost-geodesic rays, nor that the twisting orbits are such rays. The full reverse boundary, topology, and Busemann classification remain beyond the supplied result. Its precise finite-type/boundary hypotheses and a complete source for the asymptotic need verification.
\subsection{Short English statement}
For the reversed metric, the supplied Dehn-twist asymptotic gives explicit normalized logarithmic length functions in the horofunction boundary. They distinguish unoriented curve classes and transform equivariantly after renormalizing at the basepoint. This constructs a subset of boundary points; it does not classify the whole boundary or establish Busemann realization, and the appended bordered-surface context needs clarification.
\subsection{Sources}
\begin{itemize}
\item Local registry: permanent catalogue entry, exact source code catalogue entry; \texttt{aim-topology-notes.json}, \texttt{notes[195]}.
\item AIM workshop \emph{Lipschitz metric on Teichmueller space}, Problem 3.5 attributed to C. Walsh, \url{https://aimath.org/pastworkshops/lipschitzteichproblems.pdf}. No proposal year, original remarks, or original literature reference is supplied.
\item The extracted record appends a paragraph on finite-type bordered reduced Teichm\"uller space; the verification field identifies source index 195 but does not give a checked separation of that paragraph from the question.
\item Supplied research attempt 1, registry updated 14 August 2026: partial result, open-at-run label, low novelty confidence, machine-generated and not peer reviewed. It supplies the Dehn-twist asymptotic and references Walsh's proposed simple-stump Busemann formula without a title, theorem number, or complete authoritative citation.
\end{itemize}
\end{document}
